\documentclass[12pt]{article}
\usepackage[margin=1in]{geometry}
\usepackage{amsmath, amssymb}
\usepackage{booktabs}
\usepackage{hyperref}
\usepackage{natbib}
\usepackage{xcolor}
\usepackage{enumitem}

\title{Parameter Notes: QIZ Heterogeneous-Firm Trade Model\\
\large Definitions, Estimation, and Role in the Model}
\author{}
\date{\today}

\begin{document}
\maketitle
\tableofcontents
\newpage

% ─────────────────────────────────────────────────────────────────────────────
\section{Overview}
% ─────────────────────────────────────────────────────────────────────────────

This document records, parameter by parameter, what each object means
economically, how it is estimated or calibrated, and where it appears in
the model. It is updated as estimation decisions are made.

\medskip
\noindent\textbf{Sector groupings.} Two groupings are maintained throughout:
\begin{itemize}
  \item \textbf{Grouping 1 (baseline, 2 sectors):}
        $T$ = Textiles + Wearing apparel (ISIC 13--14 Rev.4 / 17--18 Rev.3);
        $O$ = All other manufacturing.
  \item \textbf{Grouping 2 (extended, 5 sectors):}
        $T$ = Textiles + Wearing apparel;
        $S_1$ = Food products (ISIC 10/15);
        $S_2$ = Chemicals (ISIC 20--21/24);
        $S_3$ = Non-metallic minerals (ISIC 23/26);
        $O$ = Residual manufacturing.
\end{itemize}
Petroleum (ISIC 19/23) and basic metals (ISIC 24/27) are absorbed into $O$
because both sectors are dominated by state-owned enterprises with very few
private exporters, making the Pareto assumption untenable for them separately.

% ─────────────────────────────────────────────────────────────────────────────
\section{Demand Parameters}
% ─────────────────────────────────────────────────────────────────────────────

\subsection{Sector demand weights $\beta_s$}

\paragraph{What it is.}
$\beta_s$ is the expenditure share of sector $s$ in the household's upper-tier
utility function. It governs how much of total manufacturing income is spent on
each sector in the baseline equilibrium.

\paragraph{Role in the model.}
Enters the upper-tier CES price index
\[
P = \Bigl(\sum_s \beta_s^{1/\eta}\, P_s^{1-1/\eta}\Bigr)^{\eta/(1-\eta)}
\]
and the sectoral expenditure allocation
\[
C_s = \beta_s \Bigl(\frac{P_s}{P}\Bigr)^{-\eta} Y,
\]
where $\eta$ is the across-sector elasticity and $Y$ is aggregate income.
A higher $\beta_s$ means sector $s$ receives a larger share of spending.

\paragraph{Estimation.}
Backed out from the Egypt 2008 Input--Output table (sheet \texttt{IO\_2})
using absorption shares:
\[
\text{Absorption}_s = \text{Output}_s + \text{Imports}_s - \text{Exports}_s,
\qquad
\beta_s = \frac{\text{Absorption}_s}{\sum_{s'} \text{Absorption}_{s'}}.
\]
Output, imports, and exports are read from columns 62, 60, and 58
respectively. Sector rows follow the ISIC Rev.4 mapping in the IO table.

\paragraph{Results.}

\medskip
\noindent\textit{Grouping 1:}
\begin{center}
\begin{tabular}{llrr}
\toprule
Sector & Description & Absorption (mn LE) & $\hat\beta_s$ \\
\midrule
$T$ & Textiles + Apparel & 34,853 & 0.0521 \\
$O$ & Other manufacturing & 634,659 & 0.9479 \\
\bottomrule
\end{tabular}
\end{center}

\noindent\textit{Grouping 2:}
\begin{center}
\begin{tabular}{llrr}
\toprule
Sector & Description & Absorption (mn LE) & $\hat\beta_s$ \\
\midrule
$T$   & Textiles + Apparel      & 34,853  & 0.0521 \\
$S_1$ & Food products           & 121,083 & 0.1809 \\
$S_2$ & Chemicals               & 61,803  & 0.0923 \\
$S_3$ & Non-metallic minerals   & 26,220  & 0.0392 \\
$O$   & Other manufacturing     & 425,553 & 0.6356 \\
\bottomrule
\end{tabular}
\end{center}

\paragraph{Script.} \texttt{estimate\_beta\_s.py}

% ─────────────────────────────────────────────────────────────────────────────
\subsection{Across-sector elasticity $\eta$}

\paragraph{What it is.}
$\eta$ governs how easily households substitute \emph{between sectors} in
response to relative price changes. It is a single number for the whole
economy.
\begin{itemize}
  \item High $\eta$: households shift spending easily across sectors when
        relative prices change.
  \item $\eta = 1$: Cobb--Douglas, expenditure shares are fixed regardless
        of prices.
  \item Low $\eta$: households maintain roughly fixed sector shares even when
        prices change a lot.
\end{itemize}

\paragraph{Role in the model.}
Enters the upper-tier utility function and aggregate price index. Determines
how a QIZ-driven price reduction in sector $T$ spills over to other sectors
through expenditure reallocation.

\paragraph{Estimation.}
Fixed from the literature. Baseline value: $\eta = 2.0$.
Robustness checks vary $\eta$.

% ─────────────────────────────────────────────────────────────────────────────
\subsection{Within-sector elasticity of substitution $\sigma_s$}

\paragraph{What it is.}
$\sigma_s$ governs how easily consumers substitute \emph{across firm-level
varieties within a sector}. It is the CES demand elasticity at the lower tier
of the utility function.
\begin{itemize}
  \item High $\sigma_s$: varieties are close substitutes, firms have little
        market power, markups are low: $\mu = \sigma_s/(\sigma_s - 1)$.
  \item Low $\sigma_s$: varieties are differentiated, firms have more market
        power, markups are high.
\end{itemize}

\paragraph{Role in the model.}
Appears in: firm revenue $r \propto \phi^{\sigma_s - 1}$; the sectoral price
index $P_s$; markups; the welfare calculation; and the upgrading gain formula
$R'/R = \delta_s^{\sigma_s - 1}$.

\paragraph{Relationship to $\xi$ and $\theta$.}
Under Pareto productivity and CES demand, the export sales distribution has
tail index $\xi_s = \theta_s / (\sigma_s - 1)$, so
$\theta_s = \xi_s \times (\sigma_s - 1)$.
$\sigma_s$ is the bridge between the observed sales distribution and the
structural productivity distribution.

\paragraph{Distinction from $\eta$.}
$\sigma_s$ governs substitution \emph{within} a sector across firm varieties
(lower tier); $\eta$ governs substitution \emph{across} sectors (upper tier).
They are completely separate objects.

\paragraph{Estimation.}
Taken from \citet{ImbsMejean2015} and the broader trade literature.
Using Broda--Weinstein estimates directly is not ideal because those identify
cross-country import substitution elasticities, not cross-firm variety
elasticities. However, for differentiated consumer goods such as
textiles and apparel, Broda--Weinstein values are a reasonable approximation.
For more homogeneous intermediate goods (metals, petroleum), Broda--Weinstein
commodity-level estimates are too low and inappropriate.

\begin{center}
\begin{tabular}{llr}
\toprule
Sector & Description & $\sigma_s$ \\
\midrule
$T$   & Textiles + Apparel    & 6.7 \\
$S_1$ & Food products         & 5.0 \\
$S_2$ & Chemicals             & 6.2 \\
$S_3$ & Non-metallic minerals & 6.0 \\
$O$   & Other manufacturing   & 6.0 \\
\bottomrule
\end{tabular}
\end{center}

% ─────────────────────────────────────────────────────────────────────────────
\section{Technology Parameters}
% ─────────────────────────────────────────────────────────────────────────────

\subsection{Labor share in production $\alpha_s$}

\paragraph{What it is.}
$\alpha_s$ is the labor cost share in variable production costs, under the
Cobb--Douglas production function
\[
y = \phi\, l^{\alpha_s}\, m^{1-\alpha_s}.
\]
It determines how intensively each sector uses labor relative to intermediate
inputs. A higher $\alpha_s$ means the sector is more labor-intensive.

\paragraph{Role in the model.}
Enters: the marginal cost function $c = (w/\alpha_s)^{\alpha_s}
(p_m/(1-\alpha_s))^{1-\alpha_s} / \phi$; the variable cost of compliance
(since Israeli sourcing changes $p_m$); the labor demand equations; and the
general equilibrium labor market clearing conditions.

\paragraph{Estimation.}
Estimated from the Egypt 2008 Input--Output table using cost shares:
\[
\hat\alpha_s =
\frac{\text{Labor Compensation}_s}
{\text{Labor Compensation}_s + \text{Intermediate Input Expenditure}_s}.
\]
Wages are in row 62 and intermediate inputs in row 57 of sheet
\texttt{IO\_2}. Column index equals row index minus 6.
Wages and intermediates are summed across constituent IO rows before taking
the ratio, so $\alpha_s$ reflects the sector-aggregate cost share.

\paragraph{Results.}

\medskip
\noindent\textit{Grouping 1:}
\begin{center}
\begin{tabular}{llrrr}
\toprule
Sector & Description & Wages (mn LE) & Intermediates (mn LE) & $\hat\alpha_s$ \\
\midrule
$T$ & Textiles + Apparel  & 4,806  & 21,101  & 0.1855 \\
$O$ & Other manufacturing & 24,689 & 360,196 & 0.0641 \\
\bottomrule
\end{tabular}
\end{center}

\noindent\textit{Grouping 2:}
\begin{center}
\begin{tabular}{llrrr}
\toprule
Sector & Description & Wages & Intermediates & $\hat\alpha_s$ \\
\midrule
$T$   & Textiles + Apparel    & 4,806  & 21,101  & 0.1855 \\
$S_1$ & Food products         & 4,647  & 75,438  & 0.0580 \\
$S_2$ & Chemicals             & 3,073  & 20,746  & 0.1290 \\
$S_3$ & Non-metallic minerals & 1,986  & 14,617  & 0.1196 \\
$O$   & Other manufacturing   & 14,984 & 249,396 & 0.0567 \\
\bottomrule
\end{tabular}
\end{center}

\noindent Note: $\alpha_T = 0.186$ is notably higher than other sectors,
consistent with textiles/apparel being a labor-intensive industry and with
the QIZ narrative that these sectors provide significant employment.
The low $\alpha_O = 0.064$ in grouping 1 reflects the dominance of
petroleum and basic metals in that residual, both of which are
intermediate-input intensive.

\paragraph{Script.} \texttt{estimate\_alpha\_s.py}

% ─────────────────────────────────────────────────────────────────────────────
\subsection{Pareto shape parameter $\theta_s$}

\paragraph{What it is.}
$\theta_s$ governs the dispersion of firm productivity within a sector.
Under the Pareto assumption, the complementary CDF of productivity satisfies
\[
\Pr(\phi > x) = \Bigl(\frac{\phi_{\min}}{x}\Bigr)^{\theta_s}.
\]
\begin{itemize}
  \item High $\theta_s$: firms are similar in productivity, distribution is
        compressed, few dominant firms.
  \item Low $\theta_s$: large productivity dispersion, fat tail of very
        productive superstar firms.
\end{itemize}

\paragraph{Distinction from $\xi$.}
$\xi_s$ is the tail index of the \emph{observed export sales distribution}.
$\theta_s$ is the tail index of the \emph{underlying productivity
distribution}. They are related by
\[
\xi_s = \frac{\theta_s}{\sigma_s - 1},
\qquad \theta_s = \xi_s \times (\sigma_s - 1).
\]
$\xi_s$ is what you measure from the data; $\theta_s$ is what the model needs.
The reason they differ is that more productive firms are amplified by
demand --- a firm twice as productive sells more than twice as much because
it charges a lower price and captures additional market share. The degree of
amplification is governed by $\sigma_s - 1$.

\paragraph{Role in the model.}
Appears in: the free entry condition (determines the equilibrium mass of
firms); export participation cutoffs (determines what share of firms export);
the aggregate price index; and the relationship between the productivity
distribution and the size distribution of exporters.

\paragraph{Requirement.}
$\theta_s > \sigma_s - 1$ must hold for the model to have a well-defined
equilibrium with finite mean firm sales. Equivalently, $\xi_s > 1$.

\paragraph{Estimation.}
$\xi_s$ is estimated from the upper tail of the firm export size distribution
in the Egyptian customs data (\texttt{final\_data\_matched.dta}), then
converted to $\theta_s$ using the literature $\sigma_s$ values above.

\medskip
\noindent\textit{Sample construction.}
\begin{itemize}
  \item Sector $T$: 2005 only, to avoid anticipation effects in the treated
        sector.
  \item All other sectors: year-by-year for 2005--2008, then $\xi$ averaged
        across years for stability.
  \item Only firm-years with \texttt{exp\_world} $> 0$ are used (active
        exporters only; filled-in panel zeros are excluded).
\end{itemize}

\noindent\textit{Estimators.} Two estimators are computed and averaged:
\begin{enumerate}
  \item \textbf{Hill (1975) MLE} with corrected threshold $x_{(k+1)}$:
        \[
        \hat\xi^{\text{Hill}} =
        \frac{1}{\frac{1}{k}\sum_{i=1}^{k} \log(x_{(i)}/x_{(k+1)})}.
        \]
  \item \textbf{OLS log-rank / log-size} with Blom (rank $- 0.5$) correction:
        \[
        \log(\text{rank}_i - 0.5) = a - \xi \cdot \log(x_i),
        \qquad \hat\xi^{\text{OLS}} = -\hat b.
        \]
\end{enumerate}
Primary estimate: $\hat\xi_s = \frac{1}{2}(\hat\xi^{\text{Hill}} +
\hat\xi^{\text{OLS}})$ at the 5\% tail cutoff (top 5\% of exporters by
value). The 5\% cutoff is preferred over 10\% or 20\% because the Hill
plot shows the Pareto slope stabilizes among the top 5\% of exporters;
including the 5--10\% range pulls in non-Pareto firms and biases $\xi$
downward.

\noindent\textit{Floor.} For sectors where $\hat\xi_s < 1$ (so that the
$\theta_s > \sigma_s - 1$ condition would be violated), $\theta_s$ is set
equal to $\sigma_s$ --- the minimum well-defined value. This applies to the
residual $O$ sector in both groupings, which contains petroleum and basic
metals (state-dominated, few private exporters, non-Pareto distribution).

\paragraph{Results.}

\medskip
\noindent\textit{Grouping 1:}
\begin{center}
\begin{tabular}{llrrrrl}
\toprule
Sector & Description & $N$ & $\sigma_s$ & $\hat\xi$ (5\%) &
$\hat\theta_s$ & Note \\
\midrule
$T$ & Textiles + Apparel  & 665   & 6.7 & 1.178 & 6.715 & estimated \\
$O$ & Other manufacturing & 15,426 & 6.0 & 0.971 & 6.000 & floored \\
\bottomrule
\end{tabular}
\end{center}

\noindent\textit{Grouping 2:}
\begin{center}
\begin{tabular}{llrrrrl}
\toprule
Sector & Description & $N$ & $\sigma_s$ & $\hat\xi$ (5\%) &
$\hat\theta_s$ & Note \\
\midrule
$T$   & Textiles + Apparel    & 665   & 6.7 & 1.178 & 6.715 & estimated \\
$S_1$ & Food products         & 2,758 & 5.0 & 1.667 & 6.666 & estimated \\
$S_2$ & Chemicals             & 2,463 & 6.2 & 1.010 & 5.254 & estimated \\
$S_3$ & Non-metallic minerals & 1,108 & 6.0 & 1.377 & 6.883 & estimated \\
$O$   & Other manufacturing   & 9,097 & 6.0 & 0.802 & 6.000 & floored \\
\bottomrule
\end{tabular}
\end{center}

\paragraph{Note on alternative estimation.}
A cleaner approach --- which removes dependence on $\sigma_s$ entirely --- is
to estimate $\theta_s$ directly from the firm TFP distribution in the 2012
census, using the production function residual
$\log\hat\phi_i = \log y_i - \hat\alpha_s \log l_i -
(1-\hat\alpha_s)\log m_i$
and applying the Hill estimator to the upper tail of $\hat\phi_i$.
This is left for future implementation.

\paragraph{Script.} \texttt{estimate\_theta\_s.py}

% ─────────────────────────────────────────────────────────────────────────────
\subsection{Lower bound of productivity $\phi_{\min,s}$}

\paragraph{What it is.}
The lower support of the Pareto productivity distribution.

\paragraph{Estimation.}
Normalized to $\phi_{\min,s} = 1$ for all sectors. This is a scale
normalization that does not affect any economically meaningful relative
productivity margin.

% ─────────────────────────────────────────────────────────────────────────────
\subsection{Productivity upgrading multiplier $\delta_s$}

\paragraph{What it is.}
$\delta_s > 1$ is the factor by which a firm's productivity scales up if it
pays the fixed upgrading cost $w_r f^U_s$. A firm with baseline productivity
$\phi$ that upgrades operates as if it had productivity $\delta_s \cdot \phi$.

\paragraph{Role in the model.}
Upgrading reduces marginal cost in all destinations simultaneously. Under
CES demand with constant markup $\mu = \sigma_s/(\sigma_s-1)$, firm revenue
in any destination $j$ satisfies $R_{rjs} \propto \phi^{\sigma_s-1}$, so
after upgrading:
\[
\frac{R'_{rjs}}{R_{rjs}} = \delta_s^{\sigma_s - 1} \quad \forall\, j.
\]
This destination-invariance is the key property: upgrading triggered by
US market access spills over to all other destinations (non-US, RoW)
proportionally. It is this spillover that identifies $\delta_s$ separately
from the direct tariff-elimination effect.

\paragraph{Estimation --- method of moments.}
The identifying assumption is that the QIZ affects non-US exports
\emph{only} through the productivity upgrading channel: tariff elimination
applies only to US exports, and Israeli input costs raise marginal costs
uniformly, so neither channel can explain a post-treatment increase in
non-US exports. Any such increase must therefore reflect upgrading.

\medskip
\noindent\textbf{Step 1: estimate the non-US export response.}
Using the matched firm panel (\texttt{final\_data\_matched.dta}, CEM weights),
we run a firm-level event study with firm fixed effects and sector$\times$year
fixed effects, baseline at $t=-1$, window $-5$ to $+6$:
\[
\log \widetilde{X}_{it}^{\mathrm{nonUS}} =
  \sum_{k \neq -1} \beta_k \cdot \mathbf{1}[t - t_i^* = k]
  + \alpha_i + \alpha_{st} + \varepsilon_{it},
\]
where $\widetilde{X}$ is winsorized at a floor of 10{,}000 USD,
$t_i^*$ is the year of first Israeli input import (treatment event),
and standard errors are clustered by firm. The post-period average
coefficient (average of $\hat\beta_0,\ldots,\hat\beta_{+6}$) is:
\[
\hat\beta^{\mathrm{post}}_{\mathrm{nonUS}} = 0.931 \quad
(SE = 0.148,\; p < 0.001).
\]
Pre-period coefficients ($t = -5$ to $-2$) are individually insignificant,
supporting the parallel trends assumption.

\medskip
\noindent\textbf{Step 2: invert the model to recover $\delta_T$.}
The model prediction is $\log R'_{RW} - \log R_{RW} = (\sigma_T - 1)\log\delta_T$,
so:
\[
\log\delta_T = \frac{\hat\beta^{\mathrm{post}}_{\mathrm{nonUS}}}{\sigma_T - 1}
= \frac{0.931}{5.7} = 0.1634,
\qquad
\hat\delta_T = e^{0.1634} = 1.178.
\]
The 95\% confidence interval, obtained by propagating the SE:
\[
\hat\delta_T \in [e^{(0.931-1.96\times0.148)/5.7},\;
                   e^{(0.931+1.96\times0.148)/5.7}]
= [1.119,\; 1.239].
\]

\paragraph{Results.}
\begin{center}
\begin{tabular}{llrr}
\toprule
Sector & Description & $\hat\delta_s$ & 95\% CI \\
\midrule
$T$ & Textiles + Apparel & 1.178 & $[1.119,\; 1.239]$ \\
$O$ & Other manufacturing & --- & not estimated \\
\bottomrule
\end{tabular}
\end{center}

\noindent $\delta_O$ is not estimated because the QIZ applies exclusively
to textile/apparel firms; non-textile firms provide the control group but
do not themselves receive treatment. The model sets $\delta_O = 1$ (no
upgrading) as a baseline, with robustness checks at $\delta_O = 1.05$.

\paragraph{Script.}
Event study: \texttt{paperv5 code.do}, Section~12.
Inversion: computed analytically from $\hat\beta^{\mathrm{post}}$ and
$\sigma_T - 1 = 5.7$.

% ─────────────────────────────────────────────────────────────────────────────
\subsection{Israeli input price premium $p_{\mathrm{IL},s}/p_{\mathrm{RW},s}$}

\paragraph{What it is.}
The ratio of the price of intermediate inputs sourced from Israel to the
price of the cheapest available alternative from the rest of the world,
for sector $s$. A value greater than 1 means Israeli inputs are more
expensive, so complying with the QIZ content requirement raises variable
production costs.

\paragraph{Role in the model.}
Enters the marginal cost of QIZ-compliant firms through the intermediate
input price $p_{m,s}$. Under Cobb--Douglas production, a firm that sources
a fraction $\gamma_s$ of its intermediates from Israel faces an effective
input price index of
\[
p_{m,s}^{\mathrm{QIZ}} = p_{\mathrm{RW},s}
  \Bigl(\frac{p_{\mathrm{IL},s}}{p_{\mathrm{RW},s}}\Bigr)^{\gamma_s}.
\]
So the compliance cost in variable terms is
$(p_{\mathrm{IL},s}/p_{\mathrm{RW},s})^{\gamma_s(1-\alpha_s)} - 1$
percent of marginal cost.

\paragraph{Estimation strategy.}
Two sources are used, with the firm survey treated as primary.

\medskip
\noindent\textbf{1. Firm survey (primary).}
A 2006 ECES survey of Egyptian QIZ exporters asked firms directly to
compare the price of Israeli inputs to their domestic/foreign substitutes.
Results:
\begin{itemize}
  \item 28 firms (60\%): Israeli inputs \textbf{more than 20\% more expensive}
  \item 11 firms (23\%): \textbf{10--20\% more expensive}
  \item 7 firms (15\%): \textbf{less than 10\% more expensive}
  \item 3 firms (6\%): Israeli inputs cheaper
\end{itemize}
The survey-implied median premium is approximately \textbf{20--25\%},
consistent with industry accounts. We use $p_{\mathrm{IL},s}/p_{\mathrm{RW},s}
= 1.20$ as the baseline estimate for all sectors (no sector breakdown
available in the survey), with robustness checks at 1.10 and 1.30.

\medskip
\noindent\textbf{2. Customs opportunity-cost estimator (corroborating).}
For each HS6 product $h$ and year $t$ in which Israel is observed as an
origin in the Egyptian customs data (2005--2008 pre-period), we compute:
\[
\text{premium}_{ht} =
\frac{uv_{\mathrm{ISR},ht} - \min_{o \neq \mathrm{ISR}} uv_{o,ht}}
     {\min_{o \neq \mathrm{ISR}} uv_{o,ht}} \times 100,
\]
where $uv$ is the USD unit value (import value divided by quantity),
restricted to observations sharing the same quantity unit within the
product-year cell (to avoid unit-of-measure confounds), and winsorized
at the 1st--99th percentile. The median of $\text{premium}_{ht}$ across
product-year cells within each sector gives the sector-level estimate.

\paragraph{Why not the log-ratio estimator?}
An earlier approach computed $\log(uv_{\mathrm{ISR}}/uv_{\mathrm{RW}})$
relative to the \emph{median} non-Israeli unit value. This produced
negative estimates (Israeli inputs appearing cheaper), inconsistent with
both the survey evidence and industry accounts. The most likely cause is
composition bias within HS6 codes: Israel tends to export simpler or
lower-grade intermediates (e.g.\ grey fabric rather than printed fabric)
within the same 6-digit code. The opportunity-cost estimator is less
sensitive to this because it uses the \emph{minimum} non-Israeli unit
value as the benchmark --- if anything biasing the premium downward ---
while the survey directly asks firms about comparable substitutes.

\paragraph{Why customs data cannot identify this parameter.}
Two approaches were attempted using Egyptian customs import microdata
(EID-Imports, 2005--2008).

\begin{enumerate}

\item \textit{Log-ratio estimator.}
For each HS6 product-year cell containing both Israeli and non-Israeli
shipments, we computed $\log(uv_{\mathrm{ISR},ht}/uv_{\mathrm{RoW},ht})$
where $uv_{\mathrm{RoW}}$ is the median non-Israeli unit value.
This produced negative estimates --- Israeli inputs appearing 15--40\%
\emph{cheaper} than the RoW median --- contradicting both the survey
evidence and industry accounts.

\item \textit{Opportunity-cost estimator.}
We replaced the RoW median with the 5th percentile of non-Israeli unit
values (the cheapest credible alternative). This produced median premiums
of over 1,000\% --- economically implausible in the opposite direction.

\end{enumerate}

Both failures reflect the same root cause: \textbf{within-HS6 composition
bias}. A 6-digit HS code typically covers a wide range of qualities, grades,
and specifications. Israel tends to export more specialized or higher-grade
intermediates within a code (e.g.\ mercerized cotton yarn vs.\ synthetic
yarn, both recorded as HS~520512). The cheapest non-Israeli shipments in the
same HS6 cell are therefore not genuine substitutes for the Israeli input.
The unit value difference measures quality mix, not the price of the same
good from different origins. No fixed-effects strategy can resolve this
without product data at a finer level than HS6, which is not available.

\paragraph{Results.}
\begin{center}
\begin{tabular}{llr}
\toprule
Source & Sector & $\hat p_{\mathrm{IL},s}/\hat p_{\mathrm{RW},s}$ \\
\midrule
Survey (ECES 2006) & All sectors & 1.20 \ (baseline) \\
Survey (ECES 2006) & All sectors & 1.10--1.30 \ (robustness range) \\
\bottomrule
\end{tabular}
\end{center}

\noindent We use $p_{\mathrm{IL},s}/p_{\mathrm{RW},s} = 1.20$ for all
sectors. No sector breakdown is available in the survey, but the composition
of Israeli imports in Q1~2006 was approximately 50\% fabrics, 15\%
chemicals, and 6\% each for zippers, thread, and cardboard, so the sample
is heavily weighted toward sector~$T$ and the 20\% figure is most directly
applicable there.

\paragraph{Script.} Survey-based; no estimation script.
Parameter stored in \texttt{params\_estimated.json} under
\texttt{price\_wedge.primary\_estimate}.

% ─────────────────────────────────────────────────────────────────────────────
\section{Policy Parameters}

\subsection{Israeli content requirement $\gamma_s$}

\paragraph{What it is.}
$\gamma_s$ is the minimum fraction of intermediate inputs that a QIZ-compliant
firm in sector $s$ must source from Israel. It enters the effective input
price index for compliant firms:
\[
p_{m,s}^{\mathrm{QIZ}} = p_{\mathrm{RW},s}
  \Bigl(\frac{p_{\mathrm{IL},s}}{p_{\mathrm{RW},s}}\Bigr)^{\gamma_s},
\]
so a higher $\gamma_s$ means a larger compliance cost when Israeli inputs
are more expensive than alternatives.

\paragraph{Estimation.}
Fixed directly from the QIZ protocol. The agreement requires that
\textbf{35\% of the value of the final product} must originate from
qualifying zones (Egypt + Israel combined), of which \textbf{at least
10.5\% must come specifically from Israel}. The relevant parameter
for the model is the Israeli-specific requirement:

\noindent\textit{Grouping 1:}
\begin{center}
\begin{tabular}{llr}
\toprule
Sector & Description & $\gamma_s$ \\
\midrule
$T$ & Textiles + Apparel & 0.105 \\
$O$ & Other manufacturing & 0.105 \\
\bottomrule
\end{tabular}
\end{center}

\noindent\textit{Grouping 2:}
\begin{center}
\begin{tabular}{llr}
\toprule
Sector & Description & $\gamma_s$ \\
\midrule
$T$   & Textiles + Apparel    & 0.105 \\
$S_1$ & Food products         & 0.105 \\
$S_2$ & Chemicals             & 0.105 \\
$S_3$ & Non-metallic minerals & 0.105 \\
$O$   & Other manufacturing   & 0.105 \\
\bottomrule
\end{tabular}
\end{center}

\noindent The 35\% total local content requirement is not a separate
parameter --- it is automatically satisfied if the Israeli 10.5\%
requirement is met, given that Egyptian value-added in QIZ firms
typically exceeds 25\%. The requirement applies uniformly across all
manufacturing sectors covered by the QIZ agreement; $\gamma_s = 0.105$
for all sectors in the multi-sector extension.

\paragraph{Robustness.}
The effective $\gamma_T$ faced by firms may differ from the statutory
10.5\% if firms source more Israeli inputs than the minimum required
(over-compliance) or if enforcement is imperfect. Robustness checks
vary $\gamma_T \in [0.08, 0.15]$.

\paragraph{Script.} No estimation script; fixed from policy documents.

% ─────────────────────────────────────────────────────────────────────────────
\section{Fixed Costs}
% ─────────────────────────────────────────────────────────────────────────────

\subsection{Destination-specific export fixed costs $f^X_{rjs}$}

\paragraph{What it is.}
$f^X_{rjs}$ is the destination-specific fixed cost (in units of labor) paid by
a firm in region $r$ and sector $s$ to serve market
$j \in \{\mathrm{US}, \mathrm{RW}\}$. It captures non-variable barriers to
exporting that are not already in the iceberg trade wedge: buyer search,
paperwork, certification, customs handling, shipping coordination, and related
market-access frictions.

\paragraph{Why this becomes a problem for QIZ uptake.}
In the strict baseline, $f^X_{Q,US,T}$ is forced to equal $f^X_{N,US,T}$ and
both regions also inherit the same Egypt-level trade wedge to the US. In that
case the model often shuts down the $(Q,T)\rightarrow \mathrm{US}$ margin or
assigns it entirely to non-QIZ firms. Once that happens the compliance choice
is never reached, because the implementation only allows US-serving QIZ-region
firms to comply. Formally, if
\[
\Pr(\mathrm{serve\ US}\mid Q,T)=0,
\]
then the model mechanically implies
\[
\Pr(\mathrm{comply}\mid Q,T)=0
\]
for any value of $f^C_T$.

\paragraph{Reduced-form Q-specific US export-access shifter.}
To generate positive QIZ take-up, the model needs an interior set of textile
firms in QIZ locations that are close to the US-export cutoff. A reduced-form
way to represent that is to allow
\[
f^X_{Q,US,T} < f^X_{N,US,T}.
\]
This should not be interpreted narrowly as lower paperwork alone. It is better
viewed as a \textbf{Q-region US-export access shifter} that summarizes better
export logistics, customs procedures, buyer matching, certification support,
and clustering of export-oriented textile firms in QIZ locations.

\paragraph{What should be done in practice.}
\begin{enumerate}
  \item Preferred option: estimate this shifter from data on Q versus non-Q
        textile firms' US export entry or US market participation.
  \item If direct estimation is not feasible, keep the common-cost model as a
        strict baseline and report the Q-specific shifter as a benchmark or
        sensitivity exercise needed to generate interior take-up.
  \item Do not present it as a literal engineering estimate of one stand-alone
        fixed cost. Present it as a reduced-form export-access wedge.
\end{enumerate}

\subsection{Compliance fixed cost $f^C_s$}

\paragraph{What it is.}
$f^C_s$ is the fixed cost (in units of labor) that a firm in the QIZ region
must pay to comply with the Israeli content requirement. It is deterministic
($\sigma_C = 0$ in the baseline), so all firms above a productivity cutoff
$\phi^C_s$ comply and all below do not.

\medskip
\noindent In the revised uptake benchmark, deterministic compliance costs are
too knife-edge: they imply either zero take-up or near-universal take-up. For
that reason the benchmark turns on compliance heterogeneity
($\sigma_C > 0$, $n_{\varepsilon} > 1$), so the model can match an interior
uptake share.

\paragraph{Target moment.}
The compliance rate among US-exporting textile firms, computed from the
matched customs panel (\texttt{final\_data\_matched.dta}):
\[
\hat p^C_T = \frac{\text{textile firms exporting to US AND importing from Israel}}
                  {\text{textile firms exporting to US}}.
\]
Year-by-year estimates (2006--2009):
\begin{center}
\begin{tabular}{lrrr}
\toprule
Sample & Denominator & Compliant & Rate \\
\midrule
Full sample 2005--2016 & 449 & 144 & \textbf{32.1\%} \\
\bottomrule
\end{tabular}
\end{center}
Denominator = textile firms that ever export to the US (2005--2016).
Numerator = of those, firms that ever import from Israel (2005--2016).
The calibration target is $\hat p^C_T = 0.32$.

\paragraph{Estimation strategy.}
$f^C_s$ cannot be solved analytically because the compliance cutoff $\phi^C_s$
depends on equilibrium wages and prices. The calibration proceeds numerically
once all other parameters are available:
\begin{enumerate}
  \item Fix all other parameters (technology, demand, trade costs, labor).
  \item For a candidate $f^C_T$, solve the full equilibrium with QIZ on.
  \item Read off \texttt{compliance\_share\_among\_US\_exporters} for $(Q, T)$.
  \item Find the $f^C_T$ that matches $\hat p^C_T = 0.32$.
\end{enumerate}
The denominator matters. The empirical moment is
\[
\Pr(\mathrm{comply}\mid \mathrm{US\ exporter}, Q,T),
\]
not compliance among all active firms. Using all active firms in the
denominator biases the target downward whenever many textile firms operate
without serving the US. Under heterogeneous compliance costs, a higher
$f^C_T$ shifts the distribution of firm-specific compliance thresholds upward
and lowers the uptake share among eligible US exporters.

\paragraph{Script.} \texttt{calibrate\_fC.py} (pending full parameter set).

% ─────────────────────────────────────────────────────────────────────────────
\subsection{Upgrading fixed cost $f^U_s$}

\paragraph{What it is.}
$f^U_s$ is the fixed cost (in units of labor) paid by a compliant firm to
scale its productivity by $\delta_s$. Under the binary upgrade mode, all
firms above a cutoff $\phi^U_s > \phi^C_s$ upgrade.

\paragraph{Target moment.}
The upgrading rate among compliant textile firms:
\[
\hat p^U_T = \frac{\text{compliant firms with }
X^{\mathrm{nonUS}}_{t+1} > X^{\mathrm{nonUS}}_{t-1}}
{\text{compliant textile firms}},
\]
where $t$ is the year of the firm's first Israeli input import
(the treatment event). From the matched customs panel
(\texttt{final\_data\_matched.dta}):

\begin{center}
\begin{tabular}{lrr}
\toprule
 & N & Rate \\
\midrule
Compliant firms with pre+post observed & 167 & --- \\
Of which non-US exports increased & 82 & \textbf{49.1\%} \\
\bottomrule
\end{tabular}
\end{center}

\noindent Denominator = textile firms that ever import from Israel (2005--2016)
with both pre- and post-treatment non-US export observations.
The calibration target is $\hat p^U_T = 0.49$.

\paragraph{Estimation strategy.}
Same sequential logic as $f^C_s$: once all other parameters are fixed,
use Brent's method to find $f^U_T$ such that the predicted upgrading rate
\texttt{upgrade\_share\_among\_active} matches $\hat p^U_T = 0.49$. The two
fixed costs are calibrated sequentially: $f^C_T$ first (compliance rate),
then $f^U_T$ conditional on $f^C_T$ (upgrading rate among compliers).

\paragraph{Script.} \texttt{compliance\_and\_upgraders.do} (moment computation).
Calibration script to be written once full parameter set is available.

% ─────────────────────────────────────────────────────────────────────────────
\subsection{Entry fixed cost $f^E_{rs}$}

\paragraph{What it is.}
$f^E_{rs}$ is the fixed cost (in units of labor) paid by a potential entrant
in region $r$ and sector $s$ before drawing its productivity from the Pareto
distribution. In equilibrium the free-entry condition pins down the mass of
active firms $M_{rs}$.

\paragraph{Role in the model.}
The ratio $f^E_{Q,s}/f^E_{N,s}$ determines the relative mass of firms
across regions for each sector. A lower entry cost in the QIZ region would
attract more entrants there. In the baseline the entry cost is assumed equal
across regions ($f^E_{Q,s} = f^E_{N,s}$) for simplicity; the firm-count
ratio then reflects differences in expected profits only.

\paragraph{Estimation strategy.}
The target moment is the ratio of sector-$s$ firm counts across regions:
\[
\frac{M_{Q,s}}{M_{N,s}} = \frac{n_{Q,s}}{n_{N,s}},
\]
taken from the pre-QIZ Egyptian Industrial Statistics (2004), which report
establishment counts by 2-digit industry and governorate.

\medskip
\noindent\textbf{Data source.}
\texttt{Data Egypt Industrial Stats 2004 (2).xlsx}, Sheet2. Contains
establishment counts (\texttt{numestabl}), employment, wages, and output
by industry $\times$ governorate for 2004 (pre-QIZ baseline).

\medskip
\noindent\textbf{QIZ region definition.}
The QIZ treatment was sub-governorate (specific \textit{qisms}), so
any governorate-level classification is an approximation. The ELMPS
treatment coding identifies the following governorates as eventually
receiving QIZ zones (census gov codes):
\begin{itemize}
  \item \textbf{2004}: Cairo (1), Giza (14), Ismailia (13), Alexandria (2),
        Port Said (3), Kafr Sheikh (21)
  \item \textbf{2005}: Fayoum (16), Beheira (12), Menya (17), Monofiyyah (11),
        Kalyoubia (19), Suez (4) [+ additional qisms in Cairo/Giza/Kafr Sheikh]
  \item \textbf{2009}: Gharbiyyah (22), Matrouh (24)
  \item \textbf{2013}: Ismailia (15), Giza qisms (18)
\end{itemize}

\medskip
\noindent\textbf{QIZ region definition.}
The 15 QIZ governorates are taken from the official IFTA Federal Register
notices and confirmed against the ELMPS treatment coding:
\begin{itemize}
  \item \textbf{2004}: Alexandria, Cairo, Port Said
  \item \textbf{2005}: Giza, Kalyoubia, Dakhliaya, Demiatte, Gharbiyyah,
        Monofiyyah, Ismailiyya, Suez
  \item \textbf{2009}: Beni Suef, Menya
  \item \textbf{2013}: Kafr Sheikh, Beheira (liberalized)
\end{itemize}
Non-QIZ (12 governorates): Sharkiah, Fayoum, Asyout, Souhag, Qena, Aswan,
Luxor, Red Sea, Wadi Jadid, Matrouh, North Sinai, South Sinai.
The full list is documented in
\texttt{Excel Sheet Data/QIZ Governorates by Treatment Year.xlsx}.

\paragraph{Results.}

\medskip
\noindent\textit{Grouping 1:}
\begin{center}
\begin{tabular}{lrrrr}
\toprule
Sector & $n_N$ & $n_Q$ & Ratio $n_Q/n_N$ \\
\midrule
$T$ & 127  & 1,013 & 7.976 \\
$O$ & 1,490 & 6,218 & 4.173 \\
\bottomrule
\end{tabular}
\end{center}

\noindent\textit{Grouping 2:}
\begin{center}
\begin{tabular}{lrrrr}
\toprule
Sector & $n_N$ & $n_Q$ & Ratio $n_Q/n_N$ \\
\midrule
$T$   & 127   & 1,013 & 7.976 \\
$S_1$ & 937   & 3,609 & 3.852 \\
$S_2$ & 124   & 534   & 4.306 \\
$S_3$ & 184   & 1,057 & 5.745 \\
$O$   & 245   & 1,018 & 4.155 \\
\bottomrule
\end{tabular}
\end{center}

\noindent Textiles has the highest ratio (7.98), consistent with QIZ zones
being concentrated in Egypt's main textile and apparel clusters
(Greater Cairo, Alexandria, Middle Delta, Suez Canal). Non-QIZ governorates
are predominantly Upper Egypt, Sinai, and frontier regions with little
manufacturing activity.

\paragraph{Script.}
Python analysis from \texttt{Data Egypt Industrial Stats 2004 (2).xlsx}
(Sheet2). Saved to \texttt{params\_estimated.json} under
\texttt{firm\_count\_ratio}.

\paragraph{Practical note.}
The textile QIZ-region share is very high in the raw governorate-level data.
In the current model this moment should be treated as a \emph{soft reference}
rather than a knife-edge target, because forcing it exactly can collapse the
interior textile US-export margin needed to identify QIZ take-up.

\subsection{Revised Full-Model Benchmark: QIZ On versus QIZ Off}

\paragraph{Purpose.}
The revised benchmark is designed to generate an interior textile uptake margin
while preserving the broad structure of the heterogeneous-firm model. It is a
benchmark for interpretation and sensitivity analysis, not yet the final
preferred baseline.

\paragraph{Key parameter changes.}
\begin{center}
\begin{tabular}{llr}
\toprule
Object & Interpretation & Value \\
\midrule
ROO cost formula & Normalized Cobb--Douglas & normalized \\
$n_{\phi}$ & Productivity grid points & 15 \\
$n_{\varepsilon}$ & Compliance-cost heterogeneity grid & 5 \\
$\sigma_C^T$ & Textile compliance-cost dispersion & 0.6 \\
$f^E_{Q,T}$ & Textile entry cost in Q region & 0.5 \\
$f^X_{Q,US,T}$ & Q-region textile US export fixed cost & 0.88 \\
$f^X_{N,US,T}$ & Non-Q textile US export fixed cost & 43.98 \\
$f^C_T$ & Textile compliance fixed cost & 2.5 \\
$f^U_T$ & Textile upgrading fixed cost & 1.8 \\
\bottomrule
\end{tabular}
\end{center}

\paragraph{Results.}
\begin{center}
\begin{tabular}{lrr}
\toprule
Moment / outcome & QIZ off & QIZ on \\
\midrule
Q-region textile firm share & 0.5234 & 0.5511 \\
$(Q,T)$ active share & 0.1333 & 0.1333 \\
$(Q,T)$ US export share among active & 1.0000 & 1.0000 \\
$(Q,T)$ compliance share among active & 0.0000 & 0.3833 \\
$(Q,T)$ compliance share among US exporters & 0.0000 & 0.3833 \\
Aggregate welfare & 325{,}381{,}023.3 & 326{,}041{,}496.4 \\
Welfare gain (\%) & --- & 0.2030 \\
$w_Q$ & 1.1639 & 1.1661 \\
$w_N$ & 0.7229 & 0.7182 \\
\bottomrule
\end{tabular}
\end{center}

\noindent The implied textile QIZ uptake rate is 38.3\%, reasonably close to
the empirical target of 32.1\%. The benchmark therefore demonstrates that the
model can generate positive take-up once the textile Q-region US-export margin
is restored and uptake is measured on the correct denominator.

% ─────────────────────────────────────────────────────────────────────────────
\section{Labor Endowments}
% ─────────────────────────────────────────────────────────────────────────────

\subsection{Regional labor force $L$, $L_Q$, $L_N$}

\paragraph{What it is.}
$L_Q$ and $L_N$ are the total labor endowments of the QIZ and non-QIZ
regions respectively. They pin down the relative size of each region and
enter the labor market clearing conditions and the wage determination
equations.

\paragraph{Estimation.}
Taken from the Egypt Annual Labor Force Survey 2004 (CAPMAS), pages 50--58,
which reports labor force by governorate. The 2004 survey is the
pre-QIZ baseline year. Governorates are assigned to QIZ or non-QIZ
following the official QIZ protocol locations.

\medskip
\noindent\textit{QIZ governorates (15):}
Cairo, Alexandria, Port Said, Suez, Demiatte, Dakhliaya, Kalyoubia,
Kafr Sheikh, Gharbiyya, Monoufiyya, Beheira, Ismailiyya, Giza,
Beni Soueif, Menya.

\noindent\textit{Non-QIZ governorates (12):}
Sharkiyya, Fayoum, Asyout, Souhag, Qena, Aswan, Luxor, Red Sea,
Wadi Jadid, Matrouh, North Sinai, South Sinai.

\paragraph{Results.}
\begin{center}
\begin{tabular}{lrr}
\toprule
Region & Labor Force & Share \\
\midrule
$L_Q$ (QIZ, 15 govs)     & 15,336,500 & 73.9\% \\
$L_N$ (non-QIZ, 12 govs) &  5,424,700 & 26.1\% \\
$L$ (total)               & 20,761,200 & 100\% \\
\bottomrule
\end{tabular}
\end{center}

\noindent The ratio $L_Q/L_N = 2.827$. The QIZ region covers 74\% of Egypt's
labor force, reflecting that it includes the major urban and Delta
governorates (Cairo, Giza, Alexandria, the textile Delta).

\paragraph{Source.}
Egypt Annual Labor Force Survey 2004, CAPMAS.
File: \texttt{Labor Force by Governorate 2004.xlsx}.

% ─────────────────────────────────────────────────────────────────────────────
\section{Parameters Still to Estimate}
% ─────────────────────────────────────────────────────────────────────────────

The following parameters remain to be estimated or calibrated.

\begin{center}
\begin{tabular}{lll}
\toprule
Parameter & Method & Data \\
\midrule
$\kappa$ & Labor mobility elasticity & ELMPS panel \\
$f^E_{rs}$ & Calibrated to firm count ratios & Industrial stats 2004 \\
$f_{rjs}$ & Calibrated to exporter participation & Customs data \\
$f^X_{Q,US,T}/f^X_{N,US,T}$ & Q-specific export-access shifter & To estimate / sensitivity \\
$d_{rjs}$ & Calibrated to export intensity & Customs data \\
$f^C_T$ & Calibrated to compliance rate (0.32) & Customs data \\
$f^U_T$ & Calibrated to upgrading rate (0.49) & Customs data \\
$E_{js}, P_{js}$ & From destination trade data & --- \\
\bottomrule
\end{tabular}
\end{center}

% ─────────────────────────────────────────────────────────────────────────────
\section{Normalizations}
% ─────────────────────────────────────────────────────────────────────────────

\begin{align}
w_N &= 1 \quad \text{(numeraire wage in non-QIZ region)} \\
\phi_{\min,s} &= 1 \quad \forall s \quad \text{(scale normalization)}
\end{align}

\bibliographystyle{plainnat}
\begin{thebibliography}{9}

\bibitem[Imbs and M\'{e}jean(2015)]{ImbsMejean2015}
Imbs, J. and I.\ M\'{e}jean (2015).
\newblock Elasticity optimism.
\newblock \textit{American Economic Journal: Macroeconomics}, 7(3):43--83.

\bibitem[Hill(1975)]{Hill1975}
Hill, B.M. (1975).
\newblock A simple general approach to inference about the tail of a
distribution.
\newblock \textit{Annals of Statistics}, 3(5):1163--1174.

\bibitem[Melitz(2003)]{Melitz2003}
Melitz, M.J. (2003).
\newblock The impact of trade on intra-industry reallocations and aggregate
industry productivity.
\newblock \textit{Econometrica}, 71(6):1695--1725.

\end{thebibliography}

\end{document}
